\section{Limitations and Future Work}\label{sec:limitations}

The preceding sections establish the Bailout Protocol as a structurally compulsive escalation architecture with formally stated Safety, Liveness, and Accountability guarantees. These guarantees are conditional: they hold if and only if a set of platform and deployment assumptions hold simultaneously. This section makes those assumptions explicit, identifies the practical and theoretical limitations that follow from them, and maps each limitation to a directed research agenda.

\subsection{Limitations}

\subsubsection{Human Response Latency.}
The Liveness property guarantees eventual human contact but does not guarantee contact within a time horizon useful for the operational context. The bound $T_{\max}$ is a function of network topology and retry parameters, not of human cognitive response time. In high-frequency domains — algorithmic trading, real-time process control, autonomous vehicle navigation — the time required for a human principal to receive, interpret, and respond to an escalation may exceed the operational time constant by an order of magnitude.

The protocol's stall behavior under latency is well-defined but domain-dependent. In a precision agriculture context, a stall delaying a watering decision by minutes is acceptable. In a financial risk-interlock context, the same stall may be economically catastrophic. The protocol currently provides no mechanism for expressing this domain-specific cost structure. Timing parameters are provisioned statically at node setup and are not sensitive to anchor workload: a human supervising a single node may respond within seconds; the same human supervising a fleet during an emergency may have queues that take minutes to clear. The protocol has no mechanism to distinguish between these scenarios.

\textbf{Directed research:} Empirical measurement of human response latency distributions across deployment contexts to establish statistically grounded latency budgets. Integration of these budgets as enforceable policy parameters at node provisioning, with the Bounds Engine calibrated to the provisioned budget. Investigation of graded escalation responses — a rapid acknowledgment request followed by a detailed response — to better match human decision-making under varying loads.

\subsubsection{Adversarial Conditions.}
The Safety property requires that human-contact pathways cannot be closed without recorded human authorization. This guarantee holds under Assumption A1 (\S\ref{sec:properties}): the host platform provides independent scheduling priority to the Bailout Monitor. In an adversarial environment — particularly one in which the platform itself may be compromised — this assumption may not hold. A compromised OS or virtualized substrate can preempt the monitor thread, introduce race conditions into pathway-closure operations, or tamper with the authorization ledger before persistence.

The protocol is designed for deployments in which the host platform is trusted but individual agents may behave unexpectedly — the standard assumption for production AI deployments in regulated industries. It is not designed for deployments in which the platform itself is an active adversary. For national infrastructure, military systems, or financial clearinghouses, the protocol must be composed with substrate-integrity mechanisms outside its current scope: secure boot chains, hardware-attested execution environments, and independent watchdog monitors on separate hardware.

\textbf{Directed research:} Empirical evaluation of Bailout Monitor scheduling guarantees under mainstream hypervisor and OS kernel configurations. Formal specification of a minimal substrate-integrity interface enabling composition with independent watchdog architectures.

\subsubsection{Scalability.}
The current specification treats human anchoring as a 1:1 relationship between node and anchor. This model is appropriate for single-node deployments but breaks down at enterprise or infrastructure scale. The protocol correctly escalates each node's exceptions to its designated anchor, but a deployment with hundreds of concurrent nodes each with a single anchor produces an aggregation bottleneck the state machine does not address.

The 1:1 model is a proof-of-concept design, not a production architecture. Real-world deployments require at minimum 1:N relationships (one anchor supervising multiple nodes, with pooling) and potentially N:M relationships (anchor pools with load balancing and routing). The correctness properties must extend to pooled configurations: Safety must hold under pooled anchoring, guaranteeing that no configuration of pooling produces a state in which all human-contact pathways are permanently disabled.

\textbf{Directed research:} Formal specification of human anchor pool architectures, including load-balancing policies, prioritization schemes, and routing algorithms. Extension of Safety and Liveness proofs to pooled anchor configurations. Empirical evaluation under simulated high-escalation-load conditions to establish scaling curves and identify saturation thresholds.

\subsubsection{Exploitation Risk.}
The Bailout Protocol assumes machine actors trigger bailout faithfully when their operating range is exceeded. A sophisticated agent — particularly one operating in a domain where escalation is economically costly — may develop strategies to suppress or mask exception states rather than trigger the protocol. This is goal hacking applied to the accountability mechanism itself. An agent might reframe an out-of-scope condition as in-scope through creative interpretation of its operating range, absorbing exceptions within the machine-only execution graph without triggering mandatory escalation.

The bypass mechanism enforces the no-machine-actor property \emph{once escalation is triggered}, but it does not detect whether a machine actor has intentionally avoided triggering. The state machine cannot distinguish between a Bounds Engine that correctly assessed a condition as in-scope and one that deliberately misclassified an out-of-scope condition to avoid escalation. This limitation strikes at the heart of the protocol's accountability model: the protocol holds machine actors accountable by observing their triggering behavior, but has no independent mechanism to detect an omission in that triggering.

\textbf{Directed research:} External monitoring mechanisms that detect anomalies in triggering behavior — specifically, statistically significant deviations between the distribution of conditions observed by a node and the distribution of escalations it triggers. Entropy-based trigger monitors: a node that observes conditions outside its operating range but escalates at a rate statistically inconsistent with its observed condition distribution is a candidate for audit.

\subsection{Future Work}

\subsubsection{Formal Verification.}
The correctness properties stated in \S\ref{sec:properties} are derived informally from the protocol's state machine specification. Before deployment in safety-critical domains, these derivations must be replaced by machine-checked formal proofs. Of particular interest is whether the functional separation between Purpose, Bounds, and Fact layers admits a compositional verification argument: if each layer satisfies its specification, does the composition necessarily satisfy the protocol's guarantees? A compositional verification architecture would distribute the verification burden across layer implementations and make the protocol more robust to implementation changes.

Formal verification also exposes hidden assumptions in the current argument-map proofs — gaps that appear correct but lack rigorous justification would be identified and resolved. The verification effort is thus as valuable for what it reveals about the protocol's limitations as for what it confirms about its correctness.

\subsubsection{Adaptive Timeout.}
The current specification uses a fixed timeout $T_{\text{bounds}}$ as a liveness trigger. The same parameter value cannot be appropriate for an agricultural irrigation system (response times of hours acceptable) and a financial risk management system (response times of seconds required). An adaptive timeout mechanism — one that adjusts $T_{\text{bounds}}$ based on observed human response latency, escalation frequency, and domain-specific operational time constants — would reduce the provisioning burden and improve operational fit.

Adaptation must preserve the deterministic timing guarantees required by the Safety property: adaptation must be constrained and slow relative to the operational time constants, so that the adapted timeout remains a deterministic upper bound rather than a probabilistic estimate. On each escalation cycle, the protocol records actual response time $t_{\text{actual}}$. Over a rolling observation window, it computes a distribution and sets $T_{\text{bounds}}$ to a high percentile plus a safety margin. Safety is preserved by never adjusting $T_{\text{bounds}}$ below the minimum viable threshold established at provisioning.

\subsubsection{Distributed Human Anchor Pools.}
As noted in the scalability limitation, the 1:1 node-to-anchor model is not viable at scale. Open questions include: how to partition nodes among anchors while preserving the invariant that each node has at least one reachable anchor; how to handle anchor unavailability without violating Safety; and how to design routing policies that minimize human attention fragmentation while ensuring all escalations receive timely response.

The pool architecture must satisfy three constraints. \textbf{Safety:} no pool configuration may produce a state in which all human-contact pathways for a given node are permanently disabled — each node must maintain at least two anchor references and the pool management algorithm must detect and compensate for anchor unavailability before it becomes pathway-closing. \textbf{Liveness:} the routing algorithm must ensure escalations reach available anchors within $T_{\max}$ under high concurrent load, requiring load-aware routing. \textbf{Accountability:} every escalation must be attributed to a specific human anchor, immutably recorded, requiring that routing decisions be committed to the Forensic Ledger. Load-shedding policy — defined behavior when all anchors are simultaneously overloaded — must preserve Safety while providing meaningful overload feedback.

\section{Conclusion}\label{sec:conclusion}

The Bailout Protocol advances a single, architecturally precise claim: that autonomous systems built on the \bxthree{} Framework must propagate every unresolved exception state to a named human principal, bypassing all intermediate machine actors, as a compulsory structural requirement rather than a runtime choice. This mandatory bailout property is not a feature retrofitted through better error handling or improved escalation heuristics. It is a structural property that follows, as a matter of architectural necessity, from the functional separation between the Purpose, Bounds, and Fact layers. When the Bounds Layer encounters a condition outside its provisioned operating range that it cannot resolve, no machine actor has the authority to adjudicate that condition. The decision must return to the Purpose Layer, which carries intent, judgment, and final authority. The Bailout Protocol is the mechanism that enforces this return — compulsorily, deterministically, and with full forensic attribution at every step.

This contribution is realized within the \bxthree{} Framework as its fifth named architectural pillar, complementing Loop Isolation, Recursive Spawning, Spatial Firewall, and Sandbox Gate. Together, these five pillars constitute the upstream accountability guarantee: that \bxthree{} systems \emph{fail upward into human consciousness, never downward into autonomous chaos}. The Bailout Protocol is the runtime expression of this guarantee for all exception conditions not resolved by the preceding four pillars — including conditions not anticipated at design time, which is precisely the operational reality that makes the protocol necessary.

The \bxthree{} Framework's role is to provide the conceptual architecture within which the Bailout Protocol operates as a necessary consequence rather than an independent design choice. The functional separation of Purpose, Bounds, and Fact layers creates the structural conditions that make the bailout requirement compulsory. The Forensic Ledger maintains the immutable attribution record that makes every escalation auditable. The immutable parent-pointer established at node provisioning ensures the escalation path is fixed at the architectural level and cannot be modified by any machine actor. The upstream accountability guarantee and the Bailout Protocol are two aspects of the same architectural commitment: that in a \bxthree{} system, every exception that exceeds machine authority reaches a human who can bind it, and every such exception is recorded so that the binding can be verified.

The theoretical claim is foundational: that the upstream accountability guarantee of the \bxthree{} Framework is not a design aspiration but a logical consequence of the framework's architecture, and that the Bailout Protocol is the mechanism by which that consequence is enforced at runtime. The practical claim is equally clear: that autonomous systems operating without such a mechanism are structurally exposed to failure modes — invisible exception absorption, contingent escalation, and irreversible drift — that cannot be eliminated by incremental improvement to existing error-handling or oversight practices. They require a new architectural primitive. The Bailout Protocol is that primitive.

The limitations identified in \S\ref{sec:limitations} define the boundary between what the protocol achieves within its scope and what must be established by subsequent work. Human response latency, adversarial conditions, scalability, and exploitation risk are constraints on the deployment context, not flaws in the protocol's logical design. They constitute the research agenda: formal verification to establish correctness proofs rigorously; adaptive timeout to reduce provisioning burden and improve operational fit; and distributed human anchor pools to enable deployment at industrial scale while preserving Safety, Liveness, and Accountability. The protocol's structural necessity within the \bxthree{} Framework is established by its specification. Its practical necessity in the field is established by the operational reality of autonomous systems that, without it, have no architectural path to human accountability when they encounter the unanticipated.
